% ================================================================
\section{Method}\label{sec:method}

\subsection{Analytic Linearization}\label{sec:linearize}

Given $\balpha$, we construct the continuous-time Jacobians
$A_c(\balpha), B_c(\balpha)$ of the 12-dimensional
Euler-angle state
$x = (p, \phi, \theta, \psi, v, \omega)$
about the flat equilibrium~\eqref{eq:eq_map}.

The $12 \times 12$ matrix $A_c$ has the block structure
\begin{equation}\label{eq:Ac}
  A_c = \begin{pmatrix}
    0_{3\times3} & 0_{3\times3} & I_3 & 0_{3\times3} \\
    0_{3\times3} & 0_{3\times3} & 0_{3\times3} & W(\phi^*,\theta^*) \\
    0_{3\times3} & G(\balpha)   & 0_{3\times3} & 0_{3\times3} \\
    0_{3\times3} & 0_{3\times3} & 0_{3\times3} & 0_{3\times3}
  \end{pmatrix},
\end{equation}
where $G(\balpha) \in \R^{3\times3}$ is the gravity coupling matrix
(derivative of the thrust vector with respect to attitude at equilibrium)
and $W(\phi^*, \theta^*)$ is the Euler-angle kinematic map
relating body angular velocity to Euler rates.

The $12 \times 4$ input matrix $B_c$ maps the four motor
thrusts to state derivatives.
It has nonzero blocks in
the velocity rows (thrust direction $\times k_t / m$,
where $k_t$ is the per-motor thrust coefficient and
$m$ is the quadrotor mass)
and the angular velocity rows
($J^{-1} M_{\mathrm{mix}}$,
where $J \in \R^{3\times3}$ is the body-frame inertia
tensor and $M_{\mathrm{mix}} \in \R^{3\times4}$ is
the motor mixing matrix that maps individual motor
thrusts to body torques via the arm geometry).

Discretization via forward Euler yields
\begin{equation}\label{eq:disc}
  A_d = I_{12} + A_c \Delta t, \qquad
  B_d = B_c \Delta t.
\end{equation}

The entire computation is closed-form trigonometry and linear algebra,
requiring no automatic differentiation.

\subsection{Flat Template MPC}\label{sec:flat_template}

Flat Template MPC extends TinyMPC by maintaining a \emph{family}
of DARE caches indexed by $\balpha$.

\textbf{Offline.}
Select a grid $\{\balpha^{(j)}\}_{j=1}^M$
(e.g., sampled along the reference trajectory).
For each grid point:
\begin{enumerate}
  \item Compute $(A_d^{(j)}, B_d^{(j)}) = \texttt{linearise}(\balpha^{(j)})$.
  \item Solve the standard (non-augmented) discrete LQR
    to obtain the terminal cost matrix $P_{\mathrm{lqr}}^{(j)}$.
  \item Solve augmented DARE~\eqref{eq:kinf} to obtain
    $K_\infty^{(j)}, P_\infty^{(j)}, C_1^{(j)}, C_2^{(j)}$.
  \item Store the full cache
    $\{A_d, B_d, K_\infty, P_\infty, C_1, C_2\}^{(j)}$.
\end{enumerate}

\textbf{Online.}
At each control step with current flat parameter $\balpha$:
\begin{enumerate}
  \item Look up the nearest cache entry
    $j^* = \arg\min_j \|\balpha - \balpha^{(j)}\|$.
  \item Compute the flat equilibrium
    $(x^*, u^*) = \Phi(\balpha)$.
  \item Form the error state $\delta x = x - x^*$.
  \item Run $K_{\mathrm{ADMM}}$ ADMM iterations using cache $j^*$
    (identical to TinyMPC's backward/forward/slack/dual passes).
  \item Apply $u = u^* + \delta u_0$, where $\delta u_0$ is
    the first-step control from the ADMM solution.
\end{enumerate}

This preserves TinyMPC's matrix-inversion-free online structure
while adapting the linearization to the current operating point.
The offline cost is $M$ DARE solves (each ${\sim}5000$ Riccati
iterations), and the online lookup is $O(M)$ nearest-neighbor search.

\subsection{Analytic Template}\label{sec:analytic}

The Analytic Template eliminates ADMM iterations entirely,
producing a control output in ${\sim}260$ arithmetic operations.

\textbf{Offline.}
\begin{enumerate}
  \item Set $\balpha_0 = (0, 0, 0, 0)$ (hover).
  \item Solve the augmented DARE at $\balpha_0$ to obtain
    the base gain $K_0 \in \R^{4 \times 12}$.
  \item For $i = 1, \ldots, 4$, compute gain sensitivities
    via central differences:
    \begin{equation}\label{eq:dK}
      \frac{\partial K}{\partial \alpha_i} \approx
      \frac{K(\balpha_0 + \epsilon \, e_i) - K(\balpha_0 - \epsilon \, e_i)}
           {2\epsilon}.
    \end{equation}
  \item Store $K_0$ and $\{\partial_i K\}_{i=1}^4$:
    a total of $5 \times 4 \times 12 = 240$ floating-point numbers.
\end{enumerate}

\textbf{Online.}
Given state $x$ and flat parameter $\balpha$:
\begin{enumerate}
  \item $\delta \balpha = \balpha - \balpha_0$.
  \item Compute $(x^*, u^*) = \Phi(\balpha)$.
  \item Interpolate the gain (192 multiply-accumulate operations):
    \begin{equation}\label{eq:K_interp}
      K(\balpha) \approx K_0 + \sum_{i=1}^{4} \delta\alpha_i \; \frac{\partial K}{\partial \alpha_i}.
    \end{equation}
  \item Compute the control (48 MACs):
    \begin{equation}\label{eq:u_analytic}
      u = u^*(\balpha) - K(\balpha) \, (x - x^*(\balpha)).
    \end{equation}
  \item (Optional) Clip to actuator bounds.
\end{enumerate}

\textbf{Relationship to TinyMPC.}
The Analytic Template output is \emph{exactly equal} to
the first forward-pass step of TinyMPC's ADMM iteration
with zero warm-start (i.e., before any constraint information
is incorporated).
Specifically, when all ADMM dual and slack variables are initialized
to zero, the first-iteration forward pass produces
$u_0 = -K_\infty \, \delta x$,
which matches~\eqref{eq:u_analytic} since $K(\balpha) \approx K_\infty(\balpha)$
from the same augmented DARE.
Subsequent ADMM iterations incorporate constraint information
through the slack and dual updates.
Thus, the Analytic Template trades constraint handling for
extreme computational efficiency.

\textbf{Connection to gain scheduling and Riccati sensitivity.}
The control law~\eqref{eq:u_analytic} with the gain
interpolation~\eqref{eq:K_interp} is precisely a
\emph{first-order gain schedule} over the flat parameter
$\balpha$~\cite{rugh2000}.
The gain sensitivities $\partial K / \partial \alpha_i$ are instances
of the classical Riccati perturbation theory~\cite{hewer1971,konstantinov1993},
evaluated numerically via~\eqref{eq:dK}.
What makes this gain schedule unusually compact is
differential flatness: whereas a generic gain schedule over the
$n$-dimensional state would require $O(n)$ sensitivity matrices,
flatness reduces the scheduling space to $\R^4$, yielding
exactly four sensitivity matrices regardless of state dimension.
This produces the identity:
\begin{equation}\label{eq:chain}
  \underbrace{u^* - K(\balpha)\,\delta x}_{\text{Analytic Template}}
  \;=\;
  \underbrace{u^* - K_\infty(\balpha)\,\delta x}_{\text{ADMM iter.\,0}}
  \;+\; O(\epsilon^2),
\end{equation}
where the $O(\epsilon^2)$ term is the gain approximation
error from the first-order Taylor expansion.

\subsection{Barrier-Augmented Template}\label{sec:barrier}

The Analytic Template handles constraints only by
clipping---the gain $K(\balpha)$ is unaware of actuator limits.
We now show that log-barrier functions can be absorbed into
the DARE cache, yielding a zero-iteration controller that is
inherently constraint-aware.

\textbf{Interior-point formulation.}
Consider box constraints $u_{\min} \leq u \leq u_{\max}$.
The log-barrier method~\cite{boyd2004,jordana2023}
replaces these with a smooth penalty:
\begin{equation}\label{eq:barrier}
  \mathcal{B}_\mu(u) = -\mu \sum_{j=1}^{m}
    \bigl[\log(u_j - u_{\min,j})
          + \log(u_{\max,j} - u_j)\bigr],
\end{equation}
where $\mu > 0$ is the barrier parameter controlling the
penalty sharpness.

\textbf{Barrier Hessian at equilibrium.}
The key observation is that the barrier's second derivative
(Hessian) at a given operating point is a \emph{known diagonal
matrix} that depends only on the distance from each input
channel to its bounds.
At the flat equilibrium $u^*(\balpha)$, this Hessian is:
\begin{equation}\label{eq:barrier_hess}
  H_\mu(\balpha) = \mu \cdot \mathrm{diag}\!\left(
    \frac{1}{(u^*_j - u_{\min,j})^2}
    + \frac{1}{(u_{\max,j} - u^*_j)^2}
  \right).
\end{equation}
This matrix is large when $u^*(\balpha)$ is near a
constraint boundary and small when it is far from all
boundaries---exactly the structure needed to make the
controller more conservative near limits.

\textbf{State barrier Hessian.}
For state constraints $x_{\min} \leq x \leq x_{\max}$
(e.g., position bounds $|p_i| \leq L$),
the same log-barrier construction yields a state barrier Hessian
at the equilibrium $x^*(\balpha)$:
\begin{equation}\label{eq:barrier_hess_x}
  H_{x,\mu}(\balpha) = \mu \cdot \mathrm{diag}\!\left(
    \frac{1}{(x^*_i - x_{\min,i})^2}
    + \frac{1}{(x_{\max,i} - x^*_i)^2}
  \right),
\end{equation}
applied only to the constrained state dimensions
(unconstrained dimensions contribute zero).
In error coordinates where $x^*_{\mathrm{pos}} = 0$,
the equilibrium position equals the reference, so the
barrier depends on the distance from the reference to
the constraint boundary:
$[H_{x,\mu}]_{i,i} = \mu \bigl(
  1/(p^*_{\mathrm{ref},i} + L)^2 + 1/(L - p^*_{\mathrm{ref},i})^2
\bigr)$.
This is $\balpha$-dependent through the reference trajectory,
making the barrier naturally time-varying.

\textbf{Barrier-augmented DARE.}
Define the barrier-augmented costs:
\begin{align}
  R_\mu(\balpha) &= R + H_\mu(\balpha), \label{eq:R_barrier} \\
  Q_\mu(\balpha) &= Q_{\mathrm{lqr}} + H_{x,\mu}(\balpha), \label{eq:Q_barrier}
\end{align}
where $Q_{\mathrm{lqr}}$ is the terminal cost from the
standard (non-augmented) discrete LQR.
Solving the augmented DARE with both $R_\mu$ and $Q_\mu$:
\begin{equation}\label{eq:dare_barrier}
  P_\mu = Q_{\mu,\rho} + A_d^\top P_\mu \bigl(
    A_d - B_d \, G_\mu^{-1} \,
    B_d^\top P_\mu A_d \bigr),
\end{equation}
where $Q_{\mu,\rho} = Q_\mu + \rho I$ and
$G_\mu = R_\mu + \rho I + B_d^\top P_\mu B_d$,
yields the barrier-augmented gain
$K_\mu(\balpha) = G_\mu^{-1} B_d^\top P_\mu A_d$.
This gain inherently reduces authority on input channels
that are near their limits and increases the effective
state cost near constraint boundaries:
when the reference position approaches the square bound,
$[H_{x,\mu}]_{i,i}$ grows, penalizing state excursions
in that direction and making the controller more conservative
precisely where constraints are tightest.

The offline and online procedures are given
in Algorithm~\ref{alg:barrier}.

\begin{algorithm}[t]
\caption{Barrier-Augmented Template MPC}
\label{alg:barrier}
\begin{algorithmic}[1]
\REQUIRE Constraint bounds $u_{\min}, u_{\max}$,
  $x_{\min}, x_{\max}$ (optional),
  barrier parameter $\mu$
\STATE \textit{--- Offline ---}
\STATE Set $\balpha_0 = (0,0,0,0)$ (hover)
\STATE $(x^*, u^*) \gets \Phi(\balpha_0)$
  \hfill\textit{flat equilibrium}
\STATE $H_\mu \gets$ \eqref{eq:barrier_hess}
  from $u^*$ and input bounds
\STATE $H_{x,\mu} \gets$ \eqref{eq:barrier_hess_x}
  from $x^*$ and state bounds
\STATE $R_\mu \gets R + H_\mu$;
  $Q_\mu \gets Q_{\mathrm{lqr}} + H_{x,\mu}$
  \hfill\textit{barrier-augmented costs}
\STATE Solve DARE~\eqref{eq:dare_barrier}
  with $R_\mu, Q_\mu \to K_{\mu,0}$
\FOR{$i = 1, \ldots, 4$}
  \STATE $K_\mu^+ \gets$ solve DARE at
    $\balpha_0 + \epsilon \, e_i$;
    $K_\mu^- \gets$ at $\balpha_0 - \epsilon \, e_i$
  \STATE $\frac{\partial K_\mu}{\partial \alpha_i}
    \gets (K_\mu^+ - K_\mu^-) / (2\epsilon)$
    \hfill\textit{central differences}
\ENDFOR
\STATE \textit{--- Online (per control step) ---}
\STATE $\balpha \gets$ current flat parameter
\STATE $K_\mu(\balpha) \approx K_{\mu,0}
  + \sum_i \delta\alpha_i \,
  \partial_i K_\mu$
  \hfill\textit{interpolate}
\STATE $(x^*, u^*) \gets \Phi(\balpha)$
\STATE $u \gets u^* - K_\mu(\balpha) \, (x - x^*)$
  \hfill\textit{${\sim}260$ ops}
\STATE \textit{--- Optional: IPM refinement ---}
\FOR{$k = 1, \ldots, K_{\mathrm{IPM}}$}
  \STATE $H_\mu \gets$ \eqref{eq:barrier_hess}
    at current $u^{(k-1)}$
  \STATE Riccati pass with updated $R_\mu$
    $\to \delta u$
  \STATE $u^{(k)} \gets u^{(k-1)} + \eta \, \delta u$
    \hfill\textit{Newton step}
\ENDFOR
\end{algorithmic}
\end{algorithm}

\textbf{$\alpha$-dependent conservatism.}
The barrier Hessian~\eqref{eq:barrier_hess} depends on
$\balpha$ through $u^*(\balpha)$.
At hover ($\balpha = 0$), the four motors produce equal
thrust at midrange, and $H_\mu$ is moderate in all channels.
At aggressive maneuvers (large $\|a\|$), some motors approach
saturation and $H_\mu$ becomes large in those channels,
automatically increasing conservatism precisely where
constraints are tightest.
The scheduling variable $\balpha$ thus simultaneously
parameterizes the linearization, the equilibrium,
\emph{and} the constraint tightness---all through a single
four-dimensional quantity.

\textbf{Online IPM refinement.}
For tighter constraint satisfaction, $K_{\mathrm{IPM}}$
Newton steps can be performed online
(lines~14--18 of Algorithm~\ref{alg:barrier}).
The initial iterate $u^{(0)}$ is the output of line~12---the
barrier-augmented template's zero-iteration output, which is
already near-feasible by construction.
Each Newton step recomputes the barrier
Hessian~\eqref{eq:barrier_hess} at the current iterate
$u^{(k-1)}$ (not at the equilibrium) and solves one Riccati
backward pass with the updated $R_\mu$---the same
$O(n^2)$ cost per step as an ADMM iteration.
Unlike ADMM (linear convergence), the Newton--IPM iteration
converges \emph{quadratically} near the
solution~\cite{boyd2004,jordana2023}: for moderate constraint
violations (the warm start is already near-feasible),
2--3 iterations typically reduce the KKT residual
below solver tolerance.

\textbf{Extended equivalence chain.}
The five variants now form a complete hierarchy:
\begin{align}\label{eq:full_chain}
  &\text{Online Relin}
  \xrightarrow{\text{cache}}
  \text{Flat Template}
  \xrightarrow{K\to 0} \notag \\
  &\text{Barrier Templ.}
  \xrightarrow{\mu\to 0}
  \text{Analytic Templ.}
  \xrightarrow{\text{1st}}
  \text{Gain Sched.}
\end{align}
The Barrier Template sits between the Flat Template
(hard constraints via $K$ ADMM iterations) and the
Analytic Template (no constraints, clip only):
it provides soft constraint satisfaction at zero online
iterations, with optional IPM refinement for hard guarantees.
Setting $\mu = 0$ recovers the unconstrained Analytic
Template; increasing $\mu$ trades tracking aggressiveness
for constraint margin.
